\begin{recipe}{Wraps}
\label{recipe:wraps}
\kicker[Brood,Basisrecept]{Brood \& basis · basisrecept}
\index[register]{Wraps}
\index[register]{Bloem}

\meta{Rusten 20–30 minuten \dvd 5 stuks}

\lettrine{Z}{elfgemaakte} wraps zijn zo gedraaid en hebben geen gist of
rijstijd nodig. Gewoon kneden, even laten rusten en dun uitrollen in de
koekenpan bakken.

\blockrule{Ingrediënten}
\begin{ingredients}
  \ing{160 g}{tarwebloem}\index[register]{Bloem}
  \ing{100 ml}{warm water (62,5\%)}
  \ing{2 el}{zonnebloemolie (ca.\ 17\%)}
  \ing{1 tl}{zout (ca.\ 4\%)}
\end{ingredients}

\blockrule{Bereiding}
\tip{Hoe dunner je de wrap uitrolt, hoe soepeler hij wordt om te vouwen of
op te rollen.}
\begin{steps}
  \item Meng de bloem en het zout in een kom. Voeg de olie en het water toe
        en meng tot een deeg.
  \item Kneed 5–10 minuten tot een soepel, elastisch deeg. Voeg een klein
        beetje extra bloem toe als het te plakkerig is.
  \item Dek af en laat het deeg 20–30 minuten rusten.
  \item Verdeel het deeg in 5 gelijke stukken en rol ze zo dun mogelijk uit.
  \item Bak elke wrap in een droge, hete koekenpan, ongeveer 30–60 seconden
        per kant tot er lichte bruine plekjes ontstaan.
  \item Stapel de wraps op en dek af met een doek, zodat ze zacht en soepel
        blijven.
\end{steps}
\end{recipe}
